\documentclass[11pt]{article}
\usepackage[utf8]{inputenc}
\usepackage[T1]{fontenc}
\usepackage{geometry}
\usepackage{hyperref}
\usepackage{booktabs}
\usepackage{longtable}
\usepackage{amsmath}
\geometry{margin=2.5cm}
\hypersetup{colorlinks=true,linkcolor=black,urlcolor=blue,citecolor=black}

\title{SIAN: metodología para construir y analizar corpus narrativos con datos públicos}
\author{Documento de trabajo para revisión editorial}
\date{29 de julio de 2026}

\begin{document}
\maketitle

\begin{abstract}
Este documento describe una metodología local para recolectar, limpiar, clasificar y analizar narrativas públicas a partir de noticias, artículos científicos abiertos, documentos institucionales, blogs, foros y otros textos web. El objetivo no es acumular grandes volúmenes de datos, sino construir un corpus trazable y balanceado que permita estudiar cómo cambia una narrativa entre fuentes, años y comunidades semánticas. La propuesta se apoya en redes complejas, procesamiento de lenguaje natural, análisis de sentimientos, extracción de marcos narrativos y un modelo multiobjetivo de cobertura de conceptos. Se toma como antecedente metodológico el trabajo publicado en Contactos sobre inteligencia artificial, narrativa mediática, redes y análisis semántico \cite{contactos625}, así como el material suplementario local sobre interacción humano-computadora y generación recursiva de narrativas visuales \cite{frontierssupp}.
\end{abstract}

\section{Problema}

Analizar una narrativa social no equivale a contar palabras. Una palabra frecuente puede ser ruido, una noticia breve puede condensar un cambio de encuadre y una conversación pública puede mostrar tensiones que no aparecen en los artículos científicos. Por eso el sistema SIAN se diseña como una plataforma de \emph{datos inteligentes}: menos dependiente de fuerza bruta y más orientada a trazabilidad, balance de fuentes y estructura.

El caso de estudio inicial es el tatuaje como práctica cultural, estética, identitaria y sanitaria. Sin embargo, la arquitectura debe funcionar para otros tópicos. El tema puede cambiar; lo que permanece es el método: definir variantes, excluir falsos positivos, recolectar por capas, limpiar, consolidar y analizar redes narrativas.

\section{Antecedente metodológico}

El artículo de Contactos \cite{contactos625} es pertinente porque usa inteligencia artificial para estudiar narrativa mediática, evolución temporal, comunidades, sentimiento y estructuras semánticas. La lección central para este proyecto es que los textos no deben tratarse como registros aislados, sino como un sistema de relaciones entre fuentes, actores, temas y cambios de percepción.

El material suplementario proporcionado \cite{frontierssupp} añade una idea útil: la interpretación humana no desaparece. El proceso combina extracción computacional con revisión humana de significados, metáforas, actores y representaciones. En SIAN esta idea se traduce en una fase explícita de validación humana: el usuario puede excluir nodos, corregir grupos, revisar actores y decidir si un resultado tiene sentido cultural.

\section{Arquitectura de recolección}

La araña mezclada es la implementación de referencia. Esto significa que la búsqueda por fuente y la búsqueda completa usan el mismo plan; sólo cambia el alcance. Una corrida completa combina capas, años, rubros y variantes. Una corrida parcial ejecuta un subconjunto del mismo plan.

\begin{longtable}{p{3cm}p{4cm}p{6cm}}
\toprule
Capa & Fuentes & Papel narrativo\\
\midrule
Noticias & GDELT, RSS, URLs semilla & Registran eventos, polémicas, casos y cambios de encuadre público.\\
Artículos científicos & OpenAlex, Crossref, Redalyc, PDF abierto & Aportan evidencia revisada o especializada; no sustituyen la conversación social.\\
Instituciones & Gobierno, reguladores, organismos & Muestran normas, alertas, políticas públicas y lenguaje oficial.\\
Foros y blogs & Foros abiertos, blogs, RSS públicos & Capturan experiencia, opinión y conversación orgánica cuando existe acceso público.\\
Otros documentos web & Páginas abiertas clasificables & Sirven como capa residual y deben mantenerse separados para no contaminar categorías.\\
\bottomrule
\end{longtable}

\section{Procesamiento}

Cada candidato pasa por la misma secuencia lógica:

\begin{enumerate}
    \item normalización de URL, DOI, medio, año y fuente;
    \item descarga del texto cuando el acceso es público y legal;
    \item limpieza de menús, anuncios, recomendaciones y restos de plantilla;
    \item exclusión de dominios o términos que producen falsos positivos;
    \item validación geográfica cuando el estudio lo requiere;
    \item deduplicación por URL canónica, DOI o huella del título;
    \item clasificación de fuente narrativa;
    \item extracción de campos para análisis local;
    \item guardado incremental en JSON;
    \item consolidación final sin duplicados.
\end{enumerate}

La limpieza debe ejecutarse antes de construir monogramas, bigramas, trigramas, actores o redes. Si se analiza texto sin normalizar, aparecen nodos irrelevantes como nombres de plantillas, marcas comerciales, menús, idiomas o restos de navegación.

\section{Análisis narrativo}

El análisis debe producir, al menos, cinco salidas:

\begin{enumerate}
    \item distribución del corpus por año, fuente, medio, país y tipo narrativo;
    \item términos, bigramas y trigramas filtrados por palabras vacías y exclusiones editables;
    \item red semántica con nodos ponderados por frecuencia y aristas ponderadas por coaparición o relación;
    \item comunidades mediante Louvain u otro método de partición;
    \item comparación temporal de marcos: problema, actor, causa, solución y consecuencia.
\end{enumerate}

La red no debe entenderse como una prueba automática de causalidad. Es un mapa exploratorio. Su valor está en hacer visibles agrupamientos, ausencias, repeticiones y desplazamientos de sentido.

\section{Intervención humana}

La intervención humana no es un adorno. Es parte del método. El usuario debe poder:

\begin{itemize}
    \item excluir nodos irrelevantes;
    \item corregir actores;
    \item agrupar frases equivalentes;
    \item separar ruido comercial de uso cultural;
    \item revisar si una comunidad semántica tiene interpretación social;
    \item decidir si el corpus está balanceado antes de publicar resultados.
\end{itemize}

\section{Límites}

El sistema no garantiza que existan 50 noticias, 50 artículos o 50 comentarios por año. Lo que sí debe garantizar es que intenta buscarlos con una estrategia transparente, registra la cobertura obtenida y reporta brechas. Una brecha no es un error si la fuente no existe o no es legalmente accesible; sí es un error si el sistema deja de buscar sin reportarlo.

\section{Conclusión}

SIAN propone una metodología reproducible para pasar de textos dispersos a un corpus narrativo analizable. Su fortaleza no está en prometer automatización total, sino en combinar extracción local, trazabilidad, balance de fuentes, redes complejas y revisión humana. Para publicación, el punto crítico será demostrar que el corpus es suficiente y balanceado para la pregunta cultural planteada.

\begin{thebibliography}{9}
\bibitem{contactos625}
Mora-Gutiérrez, R. A. y colaboradores. Artículo en \emph{Contactos}. Disponible en: \url{https://contactos.izt.uam.mx/index.php/contactos/article/view/625}.

\bibitem{frontierssupp}
Mora-Gutiérrez, R. A. y colaboradores. \emph{Human-Computer Interaction for Recursive Prompt and Image Generation on Media Narratives of Violence in Mexico}. Material suplementario local proporcionado para el proyecto, 2026.
\end{thebibliography}

\end{document}

